% ----------------------
\subsection{Section 58 (1 August 1831)}
% ----------------------
	\label{DC58}
	
	% -----
	\subsubsection{Eber D. Howe copy}
	% -----
	
		Note the final paragraph of the JSP historical background below---I've used the JSP footnotes to add the differences with Howe's version. The JSP website doesn't have a transcription of Howe's copy.
		
	% -----
	\subsubsection{Historical Background}
	% -----
		\label{DC58-HB}
		
		% --
		\paragraph{JSP}
		% --
		
			``On 1 August 1831, JS dictated this revelation to the elders of the church who had joined him in western Missouri. Just a few days earlier, a revelation had designated Jackson County, Missouri, as the location at which to build the `City of Zion.' Upon arriving in Jackson County, however, some of the elders expressed disappointment with what they found. Oliver Cowdery, Ziba Peterson, Peter Whitmer Jr., and Frederick G. Williams had been preaching to white settlers in Independence and the vicinity since they were ejected from Indian Territory west of Missouri by February 1831. Despite their efforts, those arriving in Missouri in July found fewer than ten converts, whereas some had expected a burgeoning community of believers and perhaps a settlement that would soon be able to accommodate the migration of church members. Meanwhile, tension arose between Bishop Edward Partridge and JS. The revelation of 20 July called Partridge to manage the properties of the church and `see to all things as it shall be appointed' by God's `Laws,' with the assistance of Sidney Gilbert, who had been appointed `an agent unto the church to buy lands.' According to one observer, Partridge argued with JS about the quality of the land selected for purchase. The disagreement apparently generated hard feelings on both sides. Partridge considered JS abusive, while Sidney Rigdon accused Partridge of `having insulted the Lord's prophet.'
			
			``In the wake of this incident and faced with the daunting prospect of actually building the city of Zion, JS dictated this 1 August revelation, probably at Independence. After addressing the `unbelief \& blindness of heart' of Partridge and others, the revelation gave `further directions' for the establishment of Zion, as had been promised in the 20 July revelation. The revelation instructed JS to return to Ohio, directed those appointed to build up Zion to take the initiative in moving their families to Missouri, and encouraged the elders to look beyond the land's undeveloped condition and focus on its prophesied glory. The revelation also anticipated a major migration to Jackson County and provided information about the key roles that the bishop and the agent would play in regulating that migration by making known `from time to time' the `priveliges of the lands,' or the number of individuals that the church community could accommodate. Conferences of elders were to help control the migration as well, providing counsel on who should move.
			
			``The original manuscript of this revelation is not extant. John Whitmer later copied the version featured here into Revelation Book 1, the revelation book he was keeping in Ohio. Several copies were made in addition to Whitmer's; Partridge, for example, indicated in a 5 August letter to his wife that copies of this and other Missouri revelations were to be carried to Ohio by `our brethren' returning from Missouri, and Elizabeth Van Benthusen Gilbert later showed a copy to Levi Hancock after Hancock arrived in Jackson County. Ezra Booth also had a copy of the revelation, and it is likely that others made personal copies.
			
			``Eber D. Howe, editor of the \textit{Painesville Telegraph}, denounced JS and the church in his 1834 book \textit{Mormonism Unvailed}, which published a version of the revelation `as a specimen of the manner in which the Prophet governs and rebukes his dupes.' Howe's copy differs somewhat from Whitmer's copy in Revelation Book 1: several phrases included in Whitmer's copy do not appear in Howe's, suggesting that Howe's copy may be derived from an earlier text. Howe also dated the revelation 3 August 1831, which differs from Whitmer's date of 1 August. Whitmer's copy, however, was inscribed much earlier than Howe's copy was published, and it is not known when or from whom Howe obtained a copy of the revelation, nor is it known what textual changes may have been made to Howe's copy, whether intentional or unintentional. Because Howe's version may have been copied from an earlier text than the Whitmer version, significant differences between the versions are noted in the annotation that follows.''
			
		% --
		\paragraph{HDDC}
		% --
		
			HDDC 738--740 has a few journal entries and other sources giving more background information that doesn't appear in JSP.
		
	% -----
	\subsubsection{Versions}
	% -----
		\label{DC58-V}
		
		See the \href{https://drive.google.com/file/d/1goAvNz8jS4R8slYfMG_db55Ty2z_vmgK/view?usp=sharing}{sections comparison} document.
	
		\begin{compactdesc}
			\item[JSP] \hyperref[EMS-1831-JS-1-AUG-1831]{1 August 1831}
			\item[BC] Chapter 59, 133--139
			\item[DC 1835] Section 18, 136--140
			\item[TS] 5:5 (1 March 1844), 448--450
			\item[DC 1844] Section 18, 194--200
			\item[MS] 5:9 (February 1845), 129--131
		\end{compactdesc}

	% -----
	\subsubsection{Section 58:1} % *DC58-1 
	% -----
		\label{DC58-1}

		\hyperref[HEARKEN]{HEARKEN}, O ye elders of my church, and give ear to my word, and learn of me what I will concerning you, and also concerning this land unto which I have sent you.%
			\footnote{%
				% \label{}%
				JSP note: ``The Howe copy omits the phrase `\& also concerning this land unto which I have sent you'.''
				}

		% \begin{framed}
		% \scriptsize
			% \begin{compactdesc}
				% \item[JSP] 
				% % \item[BC] 
				% % \item[]
			% \end{compactdesc}
		% \end{framed}

	% -----
	\subsubsection{Section 58:2} % *DC58-2 
	% -----
		\label{DC58-2}

		For verily I say unto you, blessed is he that keepeth my commandments, whether in life or in death; and he that is faithful in tribulation, the reward of the same is greater in the kingdom of heaven.%
			\footnote{%
				% \label{}%
				JSP note: ``The Howe copy omits `of heaven'.''
				}

		% \begin{framed}
		% \scriptsize
			% \begin{compactdesc}
				% \item[JSP] 
				% % \item[BC] 
				% % \item[]
			% \end{compactdesc}
		% \end{framed}

	% -----
	\subsubsection{Section 58:3} % *DC58-3 
	% -----
		\label{DC58-3}

		Ye cannot behold with your natural eyes,%
			\footnote{%
				% \label{}%
				For ``natural eyes,'' see \hyperref[NATURAL-PHRASES]{here}.
				}
		for the present time, the design of your God concerning those things which shall come hereafter, and the glory which shall%
			\footnote{%
				% \label{}%
				JSP note: ``The Howe copy omits `come hereafter \& the glory which shall' from this sentence, likely a scribal error.''
				}
		follow after much tribulation.

		% \begin{framed}
		% \scriptsize
			% \begin{compactdesc}
				% \item[JSP] 
				% % \item[BC] 
				% % \item[]
			% \end{compactdesc}
		% \end{framed}

	% -----
	\subsubsection{Section 58:4} % *DC58-4 
	% -----
		\label{DC58-4}

		For after much tribulation come the blessings. Wherefore the day cometh that ye shall be crowned with much glory; the hour is not yet, but is nigh at hand.

		% \begin{framed}
		% \scriptsize
			% \begin{compactdesc}
				% \item[JSP] 
				% % \item[BC] 
				% % \item[]
			% \end{compactdesc}
		% \end{framed}

	% -----
	\subsubsection{Section 58:5} % *DC58-5 
	% -----
		\label{DC58-5}

		Remember this, which I tell you before, that you may lay it to heart, and receive that which is to follow.

		% \begin{framed}
		% \scriptsize
			% \begin{compactdesc}
				% \item[JSP] 
				% % \item[BC] 
				% % \item[]
			% \end{compactdesc}
		% \end{framed}

	% -----
	\subsubsection{Section 58:6} % *DC58-6 
	% -----
		\label{DC58-6}

		Behold, verily I say unto you, for this cause I have sent you---that you might be obedient, and that your hearts might be prepared to bear testimony of the things which are to come;

		% \begin{framed}
		% \scriptsize
			% \begin{compactdesc}
				% \item[JSP] 
				% % \item[BC] 
				% % \item[]
			% \end{compactdesc}
		% \end{framed}

	% -----
	\subsubsection{Section 58:7} % *DC58-7 
	% -----
		\label{DC58-7}

		And also that you might be honored in laying the foundation, and in bearing record of the land upon which the Zion of God shall stand;

		% \begin{framed}
		% \scriptsize
			% \begin{compactdesc}
				% \item[JSP] 
				% % \item[BC] 
				% % \item[]
			% \end{compactdesc}
		% \end{framed}

	% -----
	\subsubsection{Section 58:8} % *DC58-8 
	% -----
		\label{DC58-8}

		And also that a feast of fat things might be prepared for the poor; yea, a feast of fat things, of wine on the lees well refined,%
			\footnote{%
				% \label{}%
				To say that wine is ``on the lees'' means that the wine is still resting on the sediment (the ``lees'') that settles at the bottom of the barrel or tank after fermentation. These lees are made up primarily of dead yeast cells and other particles like grape skins, seeds, and pulp. Wine aged ``on the lees'' (\textit{sur lie}, in French) can develop a richer, creamier texture, and more complex, yeasty, or nutty flavors. So if a wine is said to be ``on the lees,'' it means it's in a stage of aging with its lees, which is a process that can add depth and character to the wine.
				}
		that the earth may know that the mouths of the prophets shall not fail;

		% \begin{framed}
		% \scriptsize
			% \begin{compactdesc}
				% \item[JSP] 
				% % \item[BC] 
				% % \item[]
			% \end{compactdesc}
		% \end{framed}

	% -----
	\subsubsection{Section 58:9} % *DC58-9 
	% -----
		\label{DC58-9}

		Yea, a supper of the house of the Lord, well prepared, unto which all nations shall be invited.

		% \begin{framed}
		% \scriptsize
			% \begin{compactdesc}
				% \item[JSP] 
				% % \item[BC] 
				% % \item[]
			% \end{compactdesc}
		% \end{framed}

	% -----
	\subsubsection{Section 58:10} % *DC58-10 
	% -----
		\label{DC58-10}

		First, the rich and the learned, the wise and the noble;

		% \begin{framed}
		% \scriptsize
			% \begin{compactdesc}
				% \item[JSP] 
				% % \item[BC] 
				% % \item[]
			% \end{compactdesc}
		% \end{framed}

	% -----
	\subsubsection{Section 58:11} % *DC58-11 
	% -----
		\label{DC58-11}

		And after that cometh the day of my power; then shall the poor, the lame, and the blind, and the deaf, come in unto the marriage of the Lamb, and partake of the supper of the Lord,%
			\footnote{%
				% \label{}%
				JSP note: ``The Howe copy omits `\& partake of the supper of the Lord'.''
				}
		prepared for the great day to come.

		% \begin{framed}
		% \scriptsize
			% \begin{compactdesc}
				% \item[JSP] 
				% % \item[BC] 
				% % \item[]
			% \end{compactdesc}
		% \end{framed}

	% -----
	\subsubsection{Section 58:12} % *DC58-12 
	% -----
		\label{DC58-12}

		Behold, I, the Lord, have spoken it.

		% \begin{framed}
		% \scriptsize
			% \begin{compactdesc}
				% \item[JSP] 
				% % \item[BC] 
				% % \item[]
			% \end{compactdesc}
		% \end{framed}

	% -----
	\subsubsection{Section 58:13} % *DC58-13 
	% -----
		\label{DC58-13}

		And that the testimony might go forth from Zion, yea, from the mouth of the city of the heritage of God---

		% \begin{framed}
		% \scriptsize
			% \begin{compactdesc}
				% \item[JSP] 
				% % \item[BC] 
				% % \item[]
			% \end{compactdesc}
		% \end{framed}

	% -----
	\subsubsection{Section 58:14} % *DC58-14 
	% -----
		\label{DC58-14}

		Yea, for this cause I have sent you hither, and have selected%
			\footnote{%
				% \label{}%
				JSP note: ``The Howe copy adds `and chosen' after `Selected'.''
				}
		my servant Edward Partridge, and have appointed unto him his mission in this land.

		% \begin{framed}
		% \scriptsize
			% \begin{compactdesc}
				% \item[JSP] 
				% % \item[BC] 
				% % \item[]
			% \end{compactdesc}
		% \end{framed}

	% -----
	\subsubsection{Section 58:15} % *DC58-15 
	% -----
		\label{DC58-15}

		But if he repent not of his sins, which are unbelief and blindness of heart, let him take heed lest he fall.

		% \begin{framed}
		% \scriptsize
			% \begin{compactdesc}
				% \item[JSP] 
				% % \item[BC] 
				% % \item[]
			% \end{compactdesc}
		% \end{framed}

	% -----
	\subsubsection{Section 58:16} % *DC58-16 
	% -----
		\label{DC58-16}

		Behold his mission is given unto him, and it shall not be given again.

		% \begin{framed}
		% \scriptsize
			% \begin{compactdesc}
				% \item[JSP] 
				% % \item[BC] 
				% % \item[]
			% \end{compactdesc}
		% \end{framed}

	% -----
	\subsubsection{Section 58:17} % *DC58-17 
	% -----
		\label{DC58-17}

		And whoso standeth in this mission is appointed to be a judge in Israel, like as it was in ancient days, to divide the lands of the heritage of God unto his children;

		% \begin{framed}
		% \scriptsize
			% \begin{compactdesc}
				% \item[JSP] 
				% % \item[BC] 
				% % \item[]
			% \end{compactdesc}
		% \end{framed}

	% -----
	\subsubsection{Section 58:18} % *DC58-18 
	% -----
		\label{DC58-18}

		And to judge his people by the testimony of the just,%
			\footnote{%
				% \label{}%
				This is the only occurrence of ``testimony of the just'' in the scriptures; the only possible related verses I could find are \hyperref[DC76-51]{DC 76:51} and \hyperref[DC138-12]{DC 138:12}. But you also have places like \hyperref[DC76-69]{DC 76:69} which might help.
				} 
		and by the assistance of his counselors, according to the laws of the kingdom which are given by the prophets of God.

		% \begin{framed}
		% \scriptsize
			% \begin{compactdesc}
				% \item[JSP] 
				% % \item[BC] 
				% % \item[]
			% \end{compactdesc}
		% \end{framed}

	% -----
	\subsubsection{Section 58:19} % *DC58-19 
	% -----
		\label{DC58-19}

		For verily I say unto you, my law shall be kept on this land.

		% \begin{framed}
		% \scriptsize
			% \begin{compactdesc}
				% \item[JSP] 
				% % \item[BC] 
				% % \item[]
			% \end{compactdesc}
		% \end{framed}

	% -----
	\subsubsection{Section 58:20} % *DC58-20 
	% -----
		\label{DC58-20}

		Let no man think he is ruler; but let God rule him that judgeth, according to the counsel of his own will, or, in other words, him that counseleth or sitteth%
			\footnote{%
				% \label{}%
				JSP note: ``\textsc{text}: Or `siteth'. The Howe copy omits `councileth or' from this sentence and renders `seteth' as `sitteth'.''
				}
		upon the judgment seat.

		% \begin{framed}
		% \scriptsize
			% \begin{compactdesc}
				% \item[JSP] 
				% % \item[BC] 
				% % \item[]
			% \end{compactdesc}
		% \end{framed}

	% -----
	\subsubsection{Section 58:21} % *DC58-21 
	% -----
		\label{DC58-21}

		Let no man break the laws of the land, for he that keepeth the laws of God hath no need to break the laws of the land.%
			\footnote{%
				% \label{}%
				JSP note: ``The Howe copy omits `for he that keepeth the laws of God hath no need to break the laws of the land'.''
				}

		% \begin{framed}
		% \scriptsize
			% \begin{compactdesc}
				% \item[JSP] 
				% % \item[BC] 
				% % \item[]
			% \end{compactdesc}
		% \end{framed}

	% -----
	\subsubsection{Section 58:22} % *DC58-22 
	% -----
		\label{DC58-22}

		Wherefore, be subject to the powers that be, until he reigns whose right it is to reign, and subdues all enemies under his feet.

		% \begin{framed}
		% \scriptsize
			% \begin{compactdesc}
				% \item[JSP] 
				% % \item[BC] 
				% % \item[]
			% \end{compactdesc}
		% \end{framed}

	% -----
	\subsubsection{Section 58:23} % *DC58-23 
	% -----
		\label{DC58-23}

		Behold, the laws which ye have received from my hand are the laws of the church, and in this light ye shall hold them forth. Behold, here is wisdom.

		% \begin{framed}
		% \scriptsize
			% \begin{compactdesc}
				% \item[JSP] 
				% % \item[BC] 
				% % \item[]
			% \end{compactdesc}
		% \end{framed}

	% -----
	\subsubsection{Section 58:24} % *DC58-24 
	% -----
		\label{DC58-24}

		And now, as I spake concerning my servant Edward Partridge, this land is the land of his residence, and those whom he has appointed for his counselors; and also the land of the residence of him whom I have appointed to keep my storehouse;

		% \begin{framed}
		% \scriptsize
			% \begin{compactdesc}
				% \item[JSP] 
				% % \item[BC] 
				% % \item[]
			% \end{compactdesc}
		% \end{framed}

	% -----
	\subsubsection{Section 58:25} % *DC58-25 
	% -----
		\label{DC58-25}

		Wherefore, let them bring their families to this land, as they shall counsel between themselves and me.

		% \begin{framed}
		% \scriptsize
			% \begin{compactdesc}
				% \item[JSP] 
				% % \item[BC] 
				% % \item[]
			% \end{compactdesc}
		% \end{framed}

	% -----
	\subsubsection{Section 58:26} % *DC58-26 
	% -----
		\label{DC58-26}

		For behold, it is not meet that I should command in all things; for he that is compelled in all things, the same%
			\footnote{%
				% \label{}%
				JSP note: ``The Howe copy omits `the same'.''
				}
		is a slothful and not a wise servant; wherefore he receiveth no reward.%
			\footnote{%
				% \label{}%
				Why doesn't this person get a reward? Because they are not changing on the inside; they have not grown in their agency and intelligence in order to act to do good without being compelled or pushed by commandments and thunder. See also \hyperref[ALMA-32-16]{Alma 32:16} and \hyperref[MORONI-7-6]{Moroni 7:6--8ff}.
				}

		% \begin{framed}
		% \scriptsize
			% \begin{compactdesc}
				% \item[JSP] 
				% % \item[BC] 
				% % \item[]
			% \end{compactdesc}
		% \end{framed}

	% -----
	\subsubsection{Section 58:27} % *DC58-27 
	% -----
		\label{DC58-27}

		Verily I say, men should be anxiously engaged%
			\footnote{%
				% \label{}%
				Compare \hyperref[GALATIANS-4-18]{Galatians 4:18},
				}
		in a \hyperref[GOODNESS]{good} cause, and do many things of their own free will,%
			\footnote{%
				% \label{}%
				This is one of three scriptural uses of the phrase ``free will.'' The other two are at \hyperref[MOSIAH-18-28]{Mosiah 18:28} and \hyperref[DC124-69]{DC 124:69}.
				}
		and bring to pass much \hyperref[RIGHTEOUSNESS]{righteousness};%
			\footnote{%
				% \label{}%
				This verse, more than most, describes the way a Christian should be living---you should just be out there doing good things and doing righteous things. Christ lived this way, as we know from the record of his life, but it's also described this way in \hyperref[ACTS-10-38]{Acts 10:38}.
				}

		% \begin{framed}
		% \scriptsize
			% \begin{compactdesc}
				% \item[JSP] 
				% % \item[BC] 
				% % \item[]
			% \end{compactdesc}
		% \end{framed}

	% -----
	\subsubsection{Section 58:28} % *DC58-28 
	% -----
		\label{DC58-28}

		For the \hyperref[POWER]{power}%
			\footnote{%
				% \label{}%
				What ``power'' are we talking about here? The power to do good? If you read it this way, the verse ends up feeling like a proclamation of our free will or liberty, which isn't a bad thing; compare \hyperref[MORONI-7-13]{Moroni 7:13} and also \hyperref[2NEPHI-2-27]{2 Nephi 2:27}, where this power is enough to make us ``free to choose.'' This isn't the only reading of ``power,'' though; I think in the past I've also connected it to the priesthood and priesthood power through righteous action according to the laws of the gospel. 
				} 
		is in them, wherein they are \hyperref[AGENCY]{agents}%
			\footnote{%
				% \label{}%
				I've been reading the sections of the DC in order (it's 8 June 2022), and I think there's no question that this use of ``agent'' cannot be separated from other uses which are talking about the bishop or others, in a more legal or business sense (see \hyperref[DC58-49]{verse 49} and \hyperref[DC58-55]{verse 55} for example); we have to use those senses to understand what is going on in tihs verse, and probably what is going on in the BOM and elsewhere too, in places where it talks about ``acting for yourself.''
				}
		unto themselves. And inasmuch as men do \hyperref[GOODNESS]{good} they shall in nowise lose their reward.

		% \begin{framed}
		% \scriptsize
			% \begin{compactdesc}
				% \item[JSP] 
				% % \item[BC] 
				% % \item[]
			% \end{compactdesc}
		% \end{framed}

	% -----
	\subsubsection{Section 58:29} % *DC58-29 
	% -----
		\label{DC58-29}

		But he that doeth not anything until he is commanded, and receiveth a commandment with doubtful heart, and keepeth it with slothfulness, the same is damned.%
			\footnote{%
				% \label{}%
				This is a good example of where the LDS view of damnation as cessation of progress makes sense. Here, if you obey the commandments in this way, you aren't changing, and so you aren't progressing. This is also good for repenting as learning, because it makes me see that damnation is just where you are not learning or refusing to learn---and learning is progress. See things like \hyperref[DC93-36]{DC 93:36}.
				}

		% \begin{framed}
		% \scriptsize
			% \begin{compactdesc}
				% \item[JSP] 
				% % \item[BC] 
				% % \item[]
			% \end{compactdesc}
		% \end{framed}

	% -----
	\subsubsection{Section 58:30} % *DC58-30 
	% -----
		\label{DC58-30}

		Who am I that made man, saith the Lord, that will hold him guiltless that obeys not my commandments?

		% \begin{framed}
		% \scriptsize
			% \begin{compactdesc}
				% \item[JSP] 
				% % \item[BC] 
				% % \item[]
			% \end{compactdesc}
		% \end{framed}

	% -----
	\subsubsection{Section 58:31} % *DC58-31 
	% -----
		\label{DC58-31}

		Who am I, saith the Lord, that have promised%
			\footnote{%
				% \label{}%
				JSP note: ``The Howe copy has `ordained' instead of `promised'.''
				}
		and have not fulfilled?

		% \begin{framed}
		% \scriptsize
			% \begin{compactdesc}
				% \item[JSP] 
				% % \item[BC] 
				% % \item[]
			% \end{compactdesc}
		% \end{framed}

	% -----
	\subsubsection{Section 58:32} % *DC58-32 
	% -----
		\label{DC58-32}

		I command and men obey not; I revoke%
			\footnote{%
				% \label{}%
				Note the ``command'' and ``revoke'' language in \hyperref[DC56-4]{DC 56:4}.
				}
		and they receive not the blessing.

		% \begin{framed}
		% \scriptsize
			% \begin{compactdesc}
				% \item[JSP] 
				% % \item[BC] 
				% % \item[]
			% \end{compactdesc}
		% \end{framed}

	% -----
	\subsubsection{Section 58:33} % *DC58-33 
	% -----
		\label{DC58-33}

		Then they say in their hearts: This is not the work of the Lord, for his promises are not fulfilled. But wo unto such, for their reward lurketh beneath, and not from above.

		% \begin{framed}
		% \scriptsize
			% \begin{compactdesc}
				% \item[JSP] 
				% % \item[BC] 
				% % \item[]
			% \end{compactdesc}
		% \end{framed}

	% -----
	\subsubsection{Section 58:34} % *DC58-34 
	% -----
		\label{DC58-34}

		And now I give unto you further directions concerning this land.

		% \begin{framed}
		% \scriptsize
			% \begin{compactdesc}
				% \item[JSP] 
				% % \item[BC] 
				% % \item[]
			% \end{compactdesc}
		% \end{framed}

	% -----
	\subsubsection{Section 58:35} % *DC58-35 
	% -----
		\label{DC58-35}

		It is wisdom in me that my servant Martin Harris should be an example unto the church, in laying his moneys before the bishop of the church.

		% \begin{framed}
		% \scriptsize
			% \begin{compactdesc}
				% \item[JSP] 
				% % \item[BC] 
				% % \item[]
			% \end{compactdesc}
		% \end{framed}

	% -----
	\subsubsection{Section 58:36} % *DC58-36 
	% -----
		\label{DC58-36}

		And also, this is a law unto every man that cometh unto this land to receive an inheritance; and he shall do with his moneys according as the law directs.

		% \begin{framed}
		% \scriptsize
			% \begin{compactdesc}
				% \item[JSP] 
				% % \item[BC] 
				% % \item[]
			% \end{compactdesc}
		% \end{framed}

	% -----
	\subsubsection{Section 58:37} % *DC58-37 
	% -----
		\label{DC58-37}

		And it is wisdom also that there should be lands purchased in Independence, for the place of the storehouse, and also for the house of the printing.

		% \begin{framed}
		% \scriptsize
			% \begin{compactdesc}
				% \item[JSP] 
				% % \item[BC] 
				% % \item[]
			% \end{compactdesc}
		% \end{framed}

	% -----
	\subsubsection{Section 58:38} % *DC58-38 
	% -----
		\label{DC58-38}

		And other directions concerning my servant Martin Harris shall be given him%
			\footnote{%
				% \label{}%
				JSP note: ``The Howe copy omits `shall be given him'.''
				}
		of the Spirit, that he may receive his inheritance as seemeth him good;

		% \begin{framed}
		% \scriptsize
			% \begin{compactdesc}
				% \item[JSP] 
				% % \item[BC] 
				% % \item[]
			% \end{compactdesc}
		% \end{framed}

	% -----
	\subsubsection{Section 58:39} % *DC58-39 
	% -----
		\label{DC58-39}

		And let him repent of his sins, for he seeketh the praise of the world.

		% \begin{framed}
		% \scriptsize
			% \begin{compactdesc}
				% \item[JSP] 
				% % \item[BC] 
				% % \item[]
			% \end{compactdesc}
		% \end{framed}

	% -----
	\subsubsection{Section 58:40} % *DC58-40 
	% -----
		\label{DC58-40}

		And also let my servant William W. Phelps stand in the office to which I have appointed him, and receive his inheritance in the land;

		% \begin{framed}
		% \scriptsize
			% \begin{compactdesc}
				% \item[JSP] 
				% % \item[BC] 
				% % \item[]
			% \end{compactdesc}
		% \end{framed}

	% -----
	\subsubsection{Section 58:41} % *DC58-41 
	% -----
		\label{DC58-41}

		And also he hath need to repent, for I, the Lord, am not well pleased with him, for he seeketh to excel, and he is not sufficiently meek before me.

		% \begin{framed}
		% \scriptsize
			% \begin{compactdesc}
				% \item[JSP] 
				% % \item[BC] 
				% % \item[]
			% \end{compactdesc}
		% \end{framed}

	% -----
	\subsubsection{Section 58:42} % *DC58-42 
	% -----
		\label{DC58-42}

		Behold, he who has repented of his sins, the same is forgiven, and I, the Lord, remember them no more.

		% \begin{framed}
		% \scriptsize
			% \begin{compactdesc}
				% \item[JSP] 
				% % \item[BC] 
				% % \item[]
			% \end{compactdesc}
		% \end{framed}

	% -----
	\subsubsection{Section 58:43} % *DC58-43 
	% -----
		\label{DC58-43}

		By this ye may know if a man repenteth of his sins---behold, he will confess them and forsake them.

		% \begin{framed}
		% \scriptsize
			% \begin{compactdesc}
				% \item[JSP] 
				% % \item[BC] 
				% % \item[]
			% \end{compactdesc}
		% \end{framed}

	% -----
	\subsubsection{Section 58:44} % *DC58-44 
	% -----
		\label{DC58-44}

		And now, verily, I say concerning the residue of the elders of my church, the time has not yet come, for many years, for them to receive their inheritance in this land, except they desire it through the prayer of faith, only as it shall be appointed unto them of the Lord.

		% \begin{framed}
		% \scriptsize
			% \begin{compactdesc}
				% \item[JSP] 
				% % \item[BC] 
				% % \item[]
			% \end{compactdesc}
		% \end{framed}

	% -----
	\subsubsection{Section 58:45} % *DC58-45 
	% -----
		\label{DC58-45}

		For, behold, they shall push the people together from the ends of the earth.

		% \begin{framed}
		% \scriptsize
			% \begin{compactdesc}
				% \item[JSP] 
				% % \item[BC] 
				% % \item[]
			% \end{compactdesc}
		% \end{framed}

	% -----
	\subsubsection{Section 58:46} % *DC58-46 
	% -----
		\label{DC58-46}

		Wherefore, assemble yourselves together; and they who are not appointed to stay in this land, let them preach the gospel in the regions round about; and after that let them return to their homes.

		% \begin{framed}
		% \scriptsize
			% \begin{compactdesc}
				% \item[JSP] 
				% % \item[BC] 
				% % \item[]
			% \end{compactdesc}
		% \end{framed}

	% -----
	\subsubsection{Section 58:47} % *DC58-47 
	% -----
		\label{DC58-47}

		Let them preach by the way, and bear testimony of the truth in all places, and call upon the rich, the high and the low, and the poor to repent.

		% \begin{framed}
		% \scriptsize
			% \begin{compactdesc}
				% \item[JSP] 
				% % \item[BC] 
				% % \item[]
			% \end{compactdesc}
		% \end{framed}

	% -----
	\subsubsection{Section 58:48} % *DC58-48 
	% -----
		\label{DC58-48}

		And let them build up churches, inasmuch as the inhabitants of the earth will repent.

		% \begin{framed}
		% \scriptsize
			% \begin{compactdesc}
				% \item[JSP] 
				% % \item[BC] 
				% % \item[]
			% \end{compactdesc}
		% \end{framed}

	% -----
	\subsubsection{Section 58:49} % *DC58-49 
	% -----
		\label{DC58-49}

		And let there be an agent appointed by the voice of the church, unto the church in Ohio, to receive moneys to purchase lands in Zion.

		% \begin{framed}
		% \scriptsize
			% \begin{compactdesc}
				% \item[JSP] 
				% % \item[BC] 
				% % \item[]
			% \end{compactdesc}
		% \end{framed}

	% -----
	\subsubsection{Section 58:50} % *DC58-50 
	% -----
		\label{DC58-50}

		And I give unto my servant Sidney Rigdon a commandment, that he shall write a description of the land of Zion, and a statement of the will of God, as it shall be made known by the Spirit unto him;

		% \begin{framed}
		% \scriptsize
			% \begin{compactdesc}
				% \item[JSP] 
				% % \item[BC] 
				% % \item[]
			% \end{compactdesc}
		% \end{framed}

	% -----
	\subsubsection{Section 58:51} % *DC58-51 
	% -----
		\label{DC58-51}

		And an epistle and subscription, to be presented%
			\footnote{%
				% \label{}%
				JSP note: ``The Howe copy omits `to be presented'.''
				}
		unto all the churches to obtain moneys, to be put into the hands of the bishop, of himself or the agent, as seemeth him good or as he shall direct, to purchase lands for an inheritance for the children of God.

		% \begin{framed}
		% \scriptsize
			% \begin{compactdesc}
				% \item[JSP] 
				% % \item[BC] 
				% % \item[]
			% \end{compactdesc}
		% \end{framed}

	% -----
	\subsubsection{Section 58:52} % *DC58-52 
	% -----
		\label{DC58-52}

		For, behold, verily I say unto you,%
			\footnote{%
				% \label{}%
				JSP note: ``The Howe copy omits `verily I say unto you'.''
				}
		the Lord willeth that the disciples and the children of men should open their hearts, even to purchase this whole region of country, as soon as time will permit.

		% \begin{framed}
		% \scriptsize
			% \begin{compactdesc}
				% \item[JSP] 
				% % \item[BC] 
				% % \item[]
			% \end{compactdesc}
		% \end{framed}

	% -----
	\subsubsection{Section 58:53} % *DC58-53 
	% -----
		\label{DC58-53}

		Behold, here is wisdom. Let them do this%
		\footnote{%
			% \label{}%
			JSP note: ``The Howe copy omits `let them do this'.''
			} 
		lest they receive none inheritance, save it be by the shedding of blood.

		% \begin{framed}
		% \scriptsize
			% \begin{compactdesc}
				% \item[JSP] 
				% % \item[BC] 
				% % \item[]
			% \end{compactdesc}
		% \end{framed}

	% -----
	\subsubsection{Section 58:54} % *DC58-54 
	% -----
		\label{DC58-54}

		And again, inasmuch as there is land obtained, let there be workmen sent forth of all kinds unto this land, to labor for the saints of God.

		% \begin{framed}
		% \scriptsize
			% \begin{compactdesc}
				% \item[JSP] 
				% % \item[BC] 
				% % \item[]
			% \end{compactdesc}
		% \end{framed}

	% -----
	\subsubsection{Section 58:55} % *DC58-55 
	% -----
		\label{DC58-55}

		Let all these things be done in order; and let the privileges of the lands be made known from time to time, by the bishop or the agent of the church.

		% \begin{framed}
		% \scriptsize
			% \begin{compactdesc}
				% \item[JSP] 
				% % \item[BC] 
				% % \item[]
			% \end{compactdesc}
		% \end{framed}

	% -----
	\subsubsection{Section 58:56} % *DC58-56 
	% -----
		\label{DC58-56}

		And let the work of the gathering be not in haste, nor by flight; but let it be done as it shall be counseled by the elders of the church at the conferences, according to the knowledge which they receive%
			\footnote{%
				% \label{}%
				JSP note: ``The Howe copy reads `which they shall receive'.''
				}
		from time to time.

		% \begin{framed}
		% \scriptsize
			% \begin{compactdesc}
				% \item[JSP] 
				% % \item[BC] 
				% % \item[]
			% \end{compactdesc}
		% \end{framed}

	% -----
	\subsubsection{Section 58:57} % *DC58-57 
	% -----
		\label{DC58-57}

		And let my servant Sidney Rigdon consecrate and dedicate this land, and the spot for the temple, unto the Lord.

		% \begin{framed}
		% \scriptsize
			% \begin{compactdesc}
				% \item[JSP] 
				% % \item[BC] 
				% % \item[]
			% \end{compactdesc}
		% \end{framed}

	% -----
	\subsubsection{Section 58:58} % *DC58-58 
	% -----
		\label{DC58-58}

		And let a conference meeting be called; and after that let my servants Sidney Rigdon and Joseph Smith, Jun., return, and also Oliver Cowdery%
			\footnote{%
				% \label{}%
				JSP note: ``The Howe copy inserts `my servant' before `Oliver'.''
				}
		with them, to accomplish the residue of the work which I have appointed unto them in their own land, and the residue as shall be ruled by the conferences.

		% \begin{framed}
		% \scriptsize
			% \begin{compactdesc}
				% \item[JSP] 
				% % \item[BC] 
				% % \item[]
			% \end{compactdesc}
		% \end{framed}

	% -----
	\subsubsection{Section 58:59} % *DC58-59 
	% -----
		\label{DC58-59}

		And let no man return from this land except he bear record by the way, of that which he knows and most assuredly believes.

		% \begin{framed}
		% \scriptsize
			% \begin{compactdesc}
				% \item[JSP] 
				% % \item[BC] 
				% % \item[]
			% \end{compactdesc}
		% \end{framed}

	% -----
	\subsubsection{Section 58:60} % *DC58-60 
	% -----
		\label{DC58-60}

		Let that which has been bestowed upon%
			\footnote{%
				% \label{}%
				JSP note: ``The Howe copy has `given to' instead of `bestowed upon'.''
				}
		Ziba Peterson be taken from him; and let him stand as a member in the church, and labor with his own hands, with the brethren, until he is sufficiently chastened for all his sins; for he confesseth them not, and he thinketh to hide them.

		% \begin{framed}
		% \scriptsize
			% \begin{compactdesc}
				% \item[JSP] 
				% % \item[BC] 
				% % \item[]
			% \end{compactdesc}
		% \end{framed}

	% -----
	\subsubsection{Section 58:61} % *DC58-61 
	% -----
		\label{DC58-61}

		Let the residue of the elders of this church, who are coming to this land, some of whom are exceedingly blessed even above measure,%
			\footnote{%
				% \label{}%
				JSP note: ``The Howe copy omits `even above measure'.''
				}
		also hold a conference%
			\footnote{%
				% \label{}%
				JSP note: ``In the copy of this revelation in Revelation Book 1, Rigdon inserted here the line `upon this land: and let my Servent Edw. direct the conference.' Rigdon's insertion may have restored a line from the original revelation inadvertently omitted by copyist John Whitmer when his eye skipped from the end of the phrase `hold a conference' to the end of the phrase `direct the conference.' However, the Howe copy of the revelation does not include the line either. This conference was apparently held on 24 August 1831. (McLellin, Journal, 24 Aug. 1831.)''
				}
		upon this land.

		% \begin{framed}
		% \scriptsize
			% \begin{compactdesc}
				% \item[JSP] 
				% % \item[BC] 
				% % \item[]
			% \end{compactdesc}
		% \end{framed}

	% -----
	\subsubsection{Section 58:62} % *DC58-62 
	% -----
		\label{DC58-62}

		And let my servant Edward Partridge direct the conference which shall be held by them.

		% \begin{framed}
		% \scriptsize
			% \begin{compactdesc}
				% \item[JSP] 
				% % \item[BC] 
				% % \item[]
			% \end{compactdesc}
		% \end{framed}

	% -----
	\subsubsection{Section 58:63} % *DC58-63 
	% -----
		\label{DC58-63}

		And let them also return, preaching the gospel by the way, bearing record of the things which are revealed unto them.

		% \begin{framed}
		% \scriptsize
			% \begin{compactdesc}
				% \item[JSP] 
				% % \item[BC] 
				% % \item[]
			% \end{compactdesc}
		% \end{framed}

	% -----
	\subsubsection{Section 58:64} % *DC58-64 
	% -----
		\label{DC58-64}

		For, verily, the sound must go forth from this place into all the world, and unto the uttermost parts of the earth---the gospel must be preached unto every creature, with signs following them that believe.

		% \begin{framed}
		% \scriptsize
			% \begin{compactdesc}
				% \item[JSP] 
				% % \item[BC] 
				% % \item[]
			% \end{compactdesc}
		% \end{framed}

	% -----
	\subsubsection{Section 58:65} % *DC58-65 
	% -----
		\label{DC58-65}

		And behold the Son of Man cometh. Amen.

		% \begin{framed}
		% \scriptsize
			% \begin{compactdesc}
				% \item[JSP] 
				% % \item[BC] 
				% % \item[]
			% \end{compactdesc}
		% \end{framed}